
\section{Modeling of the Experimental Part}
\label{sec:model}

This chapter defines the model on which all experiments are run: the network, the stochastic process of the path costs, the switching penalty, the decision problem, the policies that are compared, the metrics, the data, and the way the implementation is validated. Only the experiments that were actually run are described, and their results are in Section~4.

\subsection{Network and Candidate Paths}

Let $G = (V, E)$ be a directed graph with an origin $o$ and a destination $d$. A single flow must be routed from $o$ to $d$. Time is discrete, $t = 0, 1, 2, \dots$ The set of candidate paths is $\mathcal{P} = \{P_1, \dots, P_k\}$, where each $P_i \subseteq E$ is a simple $o$--$d$ path. The candidate set is given (for example the $k$ shortest paths) and is not part of the decision. Let $E' = P_1 \cup \dots \cup P_k$ be the set of links used by at least one candidate, and $m = \abs{E'}$.

The overlap between two paths is measured with the Jaccard index:
\begin{equation}
  J(P_i, P_j) = \frac{\abs{P_i \cap P_j}}{\abs{P_i \cup P_j}},
  \label{eq:jaccard}
\end{equation}
so $J = 1$ for identical paths and $J = 0$ for link-disjoint paths.

\subsection{Stochastic Model of Path Costs}

Each link $e \in E'$ has a congestion state $x_e(t) \in \{0, 1\}$, where $0$ means free and $1$ means congested. The states of the links are independent two-state Markov chains. Link $e$ has the transition probabilities
\begin{equation}
  \Pr[x_e(t+1) = 1 \mid x_e(t) = 0] = p_e, \qquad
  \Pr[x_e(t+1) = 0 \mid x_e(t) = 1] = q_e.
  \label{eq:markov}
\end{equation}
The stationary probability of congestion is $\pi_e = p_e/(p_e+q_e)$. The speed at which the state changes is controlled by the second eigenvalue of the chain, $\rho_e = 1 - p_e - q_e$: a value close to $1$ means slowly changing (autocorrelated) states, and a value close to $0$ or negative means fast fluctuations. The chain is specified by $(\pi_e, \rho_e)$ through $p_e = \pi_e (1-\rho_e)$ and $q_e = (1-\pi_e)(1-\rho_e)$. These are valid probabilities if and only if $\rho_e \ge 1 - 1/\max(\pi_e, 1-\pi_e)$. For $\pi_e = 0.3$ this means $\rho_e \ge -0.43$, so the smallest persistence used in the grid experiments is $-0.4$.

In the grid experiments all links have the same $(\pi, \rho)$, and $\pi$ is kept fixed while $\rho$ is varied, so that the two effects are not mixed. In the stress test and in the traces the links have different parameters (Section~3.8).

The cost of a link is $c_e(x) = c^0 + \delta x$, where $c^0$ is the free-flow cost and $\delta > 0$ is the congestion penalty; in all experiments $c^0 = \delta = 1$. The cost of a path is the sum of its link costs:
\begin{equation}
  C_i(x) = \sum_{e \in P_i} c_e(x_e) = c^0 \abs{P_i} + \delta \sum_{e \in P_i} x_e, \qquad x \in \{0,1\}^m.
  \label{eq:pathcost}
\end{equation}
Since $c^0 = \delta = 1$, all path costs are integers. This has a consequence for threshold rules (Section~3.5): the cost difference between two paths takes only a few values, so a threshold has only a few distinct effective settings.

\subsection{Path-Switching Penalty}

When the flow moves from path $P_i$ to path $P_j$, a one-time penalty $s(i,j)$ is paid, with $s(i,i) = 0$. Two models are used, with the same maximum value $s_0$ so that they can be compared directly:
\begin{itemize}
  \item Constant penalty (M1): $s(i,j) = s_0$ for all $i \neq j$.
  \item Overlap-based penalty (M2): $s(i,j) = s_0 \cdot (1 - J(P_i, P_j))$.
\end{itemize}
In M2 the penalty is proportional to the fraction of links that change, in the spirit of the reconfiguration overhead counted in changed links by Kikuchi et al.\ \cite{kikuchi2025}. Switching between similar paths is cheap, and switching between disjoint paths costs the full $s_0$. For link-disjoint paths the two models coincide. Because $1 - J \le 1$, the mean penalty of M2 is not larger than that of M1 for the same $s_0$, so a difference between the models mixes the effect of the structure with the effect of the size of the penalty.

\subsection{Decision Problem as a Markov Decision Process}

At each time step the controller observes the current path $i$ and the congestion vector $x$, and chooses the path $j$ for the current step ($j = i$ means keep). The MDP is defined by:
\begin{itemize}
  \item State: $z = (i, x) \in S = \{1, \dots, k\} \times \{0,1\}^m$, so $\abs{S} = k \cdot 2^m$.
  \item Action: $j \in \{1, \dots, k\}$.
  \item Cost: $c(z, j) = s(i,j) + C_j(x)$.
  \item Transition: $(i,x) \to (j,x')$ with probability $\Pr(x,x') = \prod_{e \in E'} \Pr(x_e,x_e')$, with the link probabilities of \eqref{eq:markov}. The congestion process does not depend on the action.
  \item Criterion: expected discounted cost with discount factor $\gamma \in (0,1)$; in all experiments $\gamma = 0.95$.
\end{itemize}
The optimal value function satisfies the Bellman equation:
\begin{equation}
  V^*(i,x) = \min_{j} \Bigl[\, s(i,j) + C_j(x) + \gamma \sum_{x'} \Pr(x,x')\, V^*(j,x') \Bigr].
  \label{eq:bellman}
\end{equation}
The operator on the right-hand side is a $\gamma$-contraction in the sup norm, so value iteration converges to the unique solution $V^*$~\cite{puterman1994,bertsekas}. The iteration stops when $\norm{V_{n+1} - V_n}_\infty < \eta (1-\gamma)/(2\gamma)$, which guarantees that the resulting policy is $\eta$-optimal; the experiments use $\eta = 10^{-9}$. The optimal policy is the minimizing action in \eqref{eq:bellman}, and ties are broken in favour of keeping the current path.

The initial state is drawn from a distribution $\mu$: the initial path $i_0$ is the path with the fewest links (ties: lowest index), and $x$ follows the stationary distribution of the chains. The performance of a policy with value function $V$ is $\Phi(\mu) = \sum_{z} \mu(z)\, V(z)$.

\subsection{Policies Compared}

For a state $(i,x)$ let $j^*(i,x)$ be the currently cheapest path, $j^*(i,x) \in \argmin_j C_j(x)$. If several paths have the same minimal cost, the one with the smallest switching penalty $s(i,j)$ is chosen, so $j^* = i$ whenever the current path is among the cheapest. Let $g(i,x) = C_i(x) - C_{j^*}(x) \ge 0$ be the gain of switching. The policies are:
\begin{itemize}
  \item Optimal (OPT): the minimizer of \eqref{eq:bellman}, computed by value iteration. It is the benchmark.
  \item Never-switch (NS): always keep the current path.
  \item Greedy (GR): always move to $j^*$, ignoring the penalty.
  \item Static threshold (TH, parameter $\theta$): switch to $j^*$ if $g > \theta$, otherwise keep. It represents the rules used in practice~\cite{kikuchi2025,kabbani2014}.
  \item Overlap-aware threshold (OA, parameters $\theta_0$, $\alpha$): switch to $j^*$ if $g > \theta_0 + \alpha \cdot s(i,j^*)$. For $\theta_0 = 0$ and $\alpha = 1$ this is the one-step lookahead rule (it pays the penalty only if it is recovered immediately), related to the least one-step cost policy of Anjali et al.~\cite{anjali2002}.
\end{itemize}
For TH and OA the parameters are tuned on a grid separately for every scenario, choosing the values that minimize $\Phi(\mu)$: $\theta \in \{0, 0.25, \dots, 3\}$, $\theta_0 \in \{0, 0.5, \dots, 2\}$ and $\alpha \in \{0, 0.5, 1, 1.5, 2\}$. This gives the best case for each family, so a bad result of a threshold policy comes from the structure of the rule and not from a poor choice of parameters.

Two remarks limit how the results can be read. First, OA with $\alpha = 0$ is TH, so the OA family contains the TH family and OA can never be worse than TH on a grid that includes these values; the comparison is between a family with two parameters and one with one. Second, because path costs are integers, the gain $g$ takes only the values $0, 1, 2, 3$ on the instances used, so TH has at most four distinct members. Rules that need a memory, such as the hysteresis of FlowBender~\cite{kabbani2014}, require a counter in the state and are not part of this study.

\subsection{Evaluation Metrics}

For a stationary policy $\sigma$, the value function is computed exactly by policy evaluation, by solving the linear system:
\begin{equation}
  V^{\sigma} = c_{\sigma} + \gamma T_{\sigma} V^{\sigma}, \qquad
  V^{\sigma} = (I - \gamma T_{\sigma})^{-1} c_{\sigma},
  \label{eq:policyeval}
\end{equation}
where $c_\sigma$ is the vector of costs under $\sigma$ and $T_\sigma$ the transition matrix of the state under $\sigma$. This avoids the noise of simulation for the main results. The metrics are:
\begin{itemize}
  \item Optimality gap: $G(\sigma) = (\Phi^{\sigma} - \Phi^*)/\Phi^*$. The main metric.
  \item Failure: a scenario is a failure of a policy family at tolerance $\varepsilon$ if the gap of its best tuned member exceeds $\varepsilon$. Results are reported for $\varepsilon = 5\%$ and $\varepsilon = 2\%$.
  \item Switching rate: expected number of switches per time step, measured by Monte Carlo simulation (1000 runs of 300 steps).
  \item Oscillation rate: fraction of switches that return to the path left at most $W = 5$ steps earlier, measured in the same simulation.
  \item Average cost per step on a trace: used in the trace-driven evaluation (E4). It is not discounted, so it is not the same criterion as $\Phi$, and its gaps are close to but not equal to those of the exact evaluation.
\end{itemize}

\subsection{Experimental Design}

The parameters of the experiments are the following; everything else is fixed ($c^0 = \delta = 1$, $\gamma = 0.95$).

\begin{table}[ht]
  \centering
  \footnotesize
  \begin{tabular}{@{}>{\raggedright\arraybackslash}p{0.22\textwidth}>{\raggedright\arraybackslash}p{0.38\textwidth}>{\raggedright\arraybackslash}p{0.28\textwidth}@{}}
    \toprule
    Parameter & Values & Purpose \\
    \midrule
    Penalty ratio $s_0/\delta$ & $0.25,\ 0.5,\ 1,\ 2,\ 4$ (grid); $0.5,\ 1,\ 2$ (stress test); $2$ (traces) & Size of the switching cost relative to the gain \\
    Persistence $\rho$ & $-0.4,\ 0,\ 0.5,\ 0.8,\ 0.95$ with $\pi = 0.3$ (grid); random per link (stress test) & Speed of the changes of the link states \\
    Penalty model & M1 (constant), M2 (overlap) & Effect of the path structure \\
    Instance & A: $k = 2$ disjoint paths ($m = 4$); B: $k = 3$ paths, two of them sharing two links ($m = 7$) & Overlap structure \\
    \bottomrule
  \end{tabular}
\end{table}

The experiments are:
\begin{itemize}
  \item \textbf{E1, baseline:} instance A, model M1, the 25 scenarios of the grid $(s_0/\delta, \rho)$. Measures the gap of the tuned TH, GR and NS.
  \item \textbf{E2, penalty and persistence sweep:} instance B, the same 25 scenarios, with M1 and with M2. Measures the gap of the tuned TH under both models and the gap of OA under M2, and draws the failure maps.
  \item \textbf{E3, stress test:} instance B, 100 random scenarios (seed 1) with a different $(\pi_e, \rho_e)$ for every link, $\pi_e \in [0.1, 0.6]$ and $\rho_e \in [-0.1, 0.95]$, and $s_0 \in \{0.5, 1, 2\}$, with M1 and M2. Tests the case where the total cost of a path does not say everything about the state of the links.
  \item \textbf{E4, trace-driven evaluation:} instance B, M2, $s_0 = 2$. The chains are estimated from training traces, the policies are tuned on the estimated chain, and they are replayed on test traces (Section~3.8).
\end{itemize}

The hypotheses to be tested are:
\textbf{H1:} for M1 the best tuned threshold stays within $\varepsilon$ of the optimum.
\textbf{H2:} the failure region of M2 contains the failure region of M1 (the same scenarios plus new ones).
\textbf{H3:} OA reduces the gap of TH below $\varepsilon$ in the scenarios where TH fails.

\subsection{Data}

The experiments use synthetic data, in three forms.
\begin{itemize}
  \item \textbf{Grid (E1, E2):} no data file. The chains are the same for all links and the quantities are computed exactly from the model.
  \item \textbf{Stress test (E3):} the parameters $(\pi_e, \rho_e)$ of every link are drawn at random in the ranges above, with a fixed seed.
  \item \textbf{Traces (E4):} the script samples 10 traces of 5000 steps of the 7 links of instance B, with a different $\pi_e$ (between $0.10$ and $0.53$) and $\rho_e$ (between $-0.03$ and $0.93$) for every link and a fixed seed. The true parameters are saved, so that the estimation can be checked. The first 5 traces are the training set and the last 5 the test set. The probabilities $p_e$ and $q_e$ are estimated from the transition counts of the training set, the policies are tuned on the estimated chain, and every policy is replayed on the test traces starting from the initial path $i_0$.
\end{itemize}
The traces are generated from the same family of models that is estimated, so they show that the pipeline works, and not that the model describes a real network. The validation on real traffic matrices (GEANT) follows the same pipeline and is described in Section~2.7.

\subsection{Validation of the Implementation}

The implementation is checked by six automatic tests:
\begin{itemize}
  \item Limit cases: for $s_0 = 0$ the optimal cost equals the greedy cost, and for a very large $s_0$ it equals the never-switch cost.
  \item Solver agreement: value iteration and policy iteration give the same value function, for both penalty models.
  \item Exact versus simulation: the Monte Carlo estimate of the cost of a threshold policy (4000 runs of 300 steps) lies within twice the half-width of its 95\% confidence interval around the exact value of \eqref{eq:policyeval}.
  \item Optimality: no threshold policy with $\theta \in \{0, 0.5, \dots, 3\}$ has a lower cost than the optimal policy.
  \item Known result: for $k = 2$ and M1, the states where the optimal policy switches have larger gains than the states where it keeps, which is the threshold structure reported in the literature~\cite{anjali2002,hamouda2020}.
\end{itemize}
The stress test (E3) is also a robustness check of the conclusions with respect to the parameters of the links. The topology is fixed (instance B), so robustness with respect to the topology is not tested.

\subsection{Tools}

The implementation uses Python with NumPy (matrix operations for value iteration and policy evaluation), Matplotlib (failure maps) and \texttt{pytest} (validation). Path enumeration and shortest paths for the real-data pipeline are coded directly (depth-first search and Dijkstra), which is enough for a graph of 23 nodes. Python was chosen because the state space is small, so the speed of the language is not a limit, and because it makes the code short and easy to check. The transition matrix of the congestion process has size $2^m \times 2^m$ and is stored densely, which limits $m$ to about ten links; this is enough for the instances used. The alternatives considered are C++ (faster, but slower to develop, and needed only for much larger networks) and MATLAB (good for matrices, but a licensed tool).